\documentclass[12pt]{article}
% T-09 wrapper: compiles STATEMENT.tex and PROOF.tex standalone.
% The two files are written for the paper (sections/theory.tex and sections/appendix_theory.tex);
% this preamble only reproduces the paper macros they need.
\usepackage[margin=1in]{geometry}
\usepackage{amsmath, amssymb, amsthm}
\usepackage{bm}
\usepackage{natbib}
\usepackage[hidelinks]{hyperref}

% paper/main.tex
\newcommand{\by}{\bm{y}}
\newcommand{\blambda}{\bm{\lambda}}
\newcommand{\R}{\mathbb{R}}
\newcommand{\E}{\mathbb{E}}
\newcommand{\norm}[1]{\left\lVert#1\right\rVert}
\newtheorem{proposition}{Proposition}
\newtheorem{theorem}{Theorem}
\newtheorem{lemma}{Lemma}
\theoremstyle{definition}
\newtheorem{assumption}{Assumption}
\theoremstyle{remark}
\newtheorem{remark}{Remark}
% paper/sections/theory.tex and appendix_theory.tex
\providecommand{\pr}{\mathrm{pr}}
\providecommand{\TV}{\mathrm{TV}}
\providecommand{\lamtv}{\lambda_{\TV}}
\providecommand{\thetas}{\theta^{*}}
\providecommand{\en}{\varepsilon_n}
\providecommand{\fl}{\underline f}
\providecommand{\fu}{\bar f}
\providecommand{\Vb}{\bar V}
\providecommand{\nn}[1]{\|#1\|_n}
\providecommand{\ind}{\mathbf 1}

\title{T-09: the joint theorem for non-crossing quantile curves}
\author{Theory department (prover)}
\date{25 November 2026}

\begin{document}
\maketitle

\noindent This file compiles the two blocks written for the paper: the statement, for the section on posterior contraction, and the proof, for the appendix of proofs. References to labels of the paper appear as ``??'' here; the list below says what each one is.

\paragraph{Cited from the paper.}
\begin{itemize}\itemsep1pt
\item \texttt{thm:main}: the contraction theorem for one level, with constants $(M,C,n_0)$ depending only on $(c,k,L,\fl,\fu,\Vb)$. Used as a black box, once per level.
\item \texttt{ass:A}, \texttt{ass:B}: Assumptions A and B. \texttt{eq:post-fixed}: the posterior of one level at $\sigma=1$. \texttt{eq:joint\_order}, \texttt{eq:joint\_target}, \texttt{subsec:joint}: the set $\mathcal C$, the joint target and the section on the joint sampler.
\item \texttt{pf:eq:post}, \texttt{pf:lem:coord}(b): the posterior in the appendix and the finiteness of its normalizing integral for $n>r$.
\item \texttt{pf:prop:shell}: named only, to say that $J$ there is the number of knots.
\item \texttt{pf:lem:evid}: the event $\Omega_E$, with $\pr(\Omega_E^c)\le\frac54(n\en^2)^{-1}$ for $n\ge n_1$.
\item \texttt{pf:lem:interp}: $\|\delta\|_\infty\le C_I\{\nn{\delta}+\nn{\delta}^{(2r-2)/(2r-1)}\TV^{(r)}(\delta)^{1/(2r-1)}\}$.
\item \texttt{pf:prop:local}: the event $\Omega_G$, with $\pr(\Omega_G^c)\le64/(\fl^2n\rho_n^2)$ for $n\ge n_2$; part (b), $\Pi_\lambda(\theta-\thetas\notin K_n\mid y)\le r!\{e^{-\fl n\rho_n^2/16}+e^{-n/8}\}$ on $\Omega_G\cap\Omega_E$; and, from the proof of part (a), $\sup_{K_n}\|\delta\|_\infty\le C_I\{\rho_n+(\rho_n^{2r-2}T_n)^{1/(2r-1)}\}$ and $\rho_n^{2r-2}T_n\to0$. Here $\rho_n^2=n^{-(2r-1)/(2r+1)}\log n$, $T_n=(4/c)\sqrt n$ and $K_n=\{\delta:\ \nn{\delta}\le\rho_n,\ \TV^{(r)}(\delta)\le T_n\}$.
\item The subsection ``Reduction to $\sigma=1$'' of the appendix (it has no label): the law of $\theta/\sigma_0$ under the $\sigma_0$-posterior is the $\sigma=1$ posterior based on $y/\sigma_0$ with $c$ replaced by $\sigma_0c$, and the rescaled data satisfy the assumptions with $(L/\sigma_0,\sigma_0\fl,\sigma_0\fu)$ and $\Vb/\sigma_0$.
\end{itemize}

\bigskip\hrule\bigskip
\noindent\textbf{Block 1: statement} (\texttt{STATEMENT.tex}).

% T-09: joint theorem, main-text block. No preamble; to be \input in Section~\ref{sec:theory}
% (suggested place: after Section~\ref{sec:main-result}). Uses the macros of sections/theory.tex.
\subsection{Non-crossing quantile curves}\label{sec:joint-result}

Theorem~\ref{thm:main} concerns a single quantile level. The joint sampler of Section~\ref{subsec:joint} fits $J\ge2$ levels $0<\tau_1<\cdots<\tau_J<1$ together and restricts the trends to the set $\mathcal C$ of \eqref{eq:joint_order}, on which they do not cross. The same rate holds at every level of the joint posterior defined below, under the assumptions of Theorem~\ref{thm:main} imposed at every level.

For $j=1,\dots,J$, let $\thetas_j=(\thetas_{1j},\dots,\thetas_{nj})$ be the vector of the $\tau_j$-quantiles of $y_1,\dots,y_n$, let $V^{*}_j=\TV(\thetas_j)$, and let $\Pi_j(\cdot\mid\by)$ be the posterior \eqref{eq:post-fixed} at the level $\tau_j$ with penalty $\lambda_j$. Write $\Theta=(\theta_1,\dots,\theta_J)$. For every Borel set $B\subset(\R^n)^J$, the joint posterior is
\begin{equation}\label{eq:joint-fixed}
\Pi_{\mathcal C}(B\mid\by)=\frac{\int_{B\cap\mathcal C}\prod_{j=1}^J\Pi_j(d\theta_j\mid\by)}{\int_{\mathcal C}\prod_{j=1}^J\Pi_j(d\theta_j\mid\by)} .
\end{equation}
This is the conditional posterior of $\Theta$ given the penalties and the scales under \eqref{eq:joint_target}, at $\sigma_1=\cdots=\sigma_J=1$. For $n>k+1$ the denominator is positive for every $\by$, so $\Pi_{\mathcal C}$ is proper (Appendix~\ref{app:proofs}). As before, $\E_{\thetas}$ is expectation under the true distribution of $\by$.

\begin{theorem}\label{thm:joint}
Let $y_1,\dots,y_n$ be independent and satisfy Assumptions~\ref{ass:A} and~\ref{ass:B} at each of the levels $\tau_1<\cdots<\tau_J$, with the same constants $L$, $\fl$ and $\fu$, and let $V^{*}_j\le\Vb$ for all $j$ and $n$. Let $\sigma_j=1$ and
\[
\lambda_j=c_j\,n^{(2k+1)/2},\qquad\text{equivalently}\qquad\lambda_{\TV,j}=c_j\,n^{1/2},\qquad j=1,\dots,J,
\]
with constants $c_j>0$. Then there exist constants $M$, $C>0$ and $n_0$, depending only on $c_1,\dots,c_J$, $k$, $L$, $\fl$, $\fu$, $\Vb$, $J$, and the smallest spacing $\min_{j<J}(\tau_{j+1}-\tau_j)$, such that for all $n\ge n_0$ and every $j=1,\dots,J$
\[
\E_{\thetas}\,\Pi_{\mathcal C}\bigl(\Theta:\ \Delta_n^2(\theta_j-\thetas_j)>M\,n^{-(2k+1)/(2k+2)}\bigm|\by\bigr)\le C\,n^{-1/(2k+2)}.
\]
The bound is uniform over all laws of $\by$ that satisfy Assumptions~\ref{ass:A} and~\ref{ass:B} at every level with the given constants and $V^{*}_j\le\Vb$. The same holds for any fixed $\sigma_j=\sigma_{0j}>0$, where $\Pi_j$ in \eqref{eq:joint-fixed} is the posterior at the scale $\sigma_{0j}$, with constants depending also on $\sigma_{01},\dots,\sigma_{0J}$.
\end{theorem}

No assumption is added beyond those of Theorem~\ref{thm:main} at each level. In particular, the true quantile curves are not assumed to be separated: Assumption~\ref{ass:B} at two fixed neighbouring levels already keeps their curves at least $\min\{L,(\tau_{j+1}-\tau_j)/\fu\}$ apart at every index. The proof, in Appendix~\ref{app:proofs}, uses this separation to show that the $J$ individual posteriors together place mass tending to one on $\mathcal C$, with probability tending to one, and then applies Theorem~\ref{thm:main} at each level.

\begin{remark}[Scope]
Theorem~\ref{thm:joint} is asymptotic. Its proof requires each individual posterior to concentrate, in the supremum norm, within half the separation $\min\{L,\min_j(\tau_{j+1}-\tau_j)/\fu\}$ guaranteed above, and the supremum-norm bound that the proof provides decreases slowly: for the default $k=2$ it is of order $n^{-13/70}$, up to a logarithmic factor. At practical sample sizes, conditioning on $\mathcal C$ can change the posterior materially, and the theorem does not describe that regime. As in Theorem~\ref{thm:main}, the penalties $\lambda_j$ and the scales $\sigma_j$ are fixed; the joint posterior in which the sampler learns them is not covered. The levels are also fixed, so a grid of levels that becomes finer as $n$ grows is not covered.
\end{remark}


\bigskip\hrule\bigskip
\noindent\textbf{Block 2: proof} (\texttt{PROOF.tex}).

% T-09: proof of the joint theorem. No preamble; to be \input in Appendix~\ref{app:proofs},
% after the proof of Theorem~\ref{thm:main}. Uses the macros of sections/appendix_theory.tex.
\subsection{Proof of Theorem~\ref{thm:joint}}\label{pf:sec:joint}

As in the rest of this appendix, $r=k+1$ and $\en^2=n^{-(2r-1)/(2r)}$. Let $\Pi^{\otimes}(\cdot\mid y)$ be the product of $\Pi_1(\cdot\mid y),\dots,\Pi_J(\cdot\mid y)$, a probability measure on $(\R^n)^J$. Then \eqref{eq:joint-fixed} reads
\[
\Pi_{\mathcal C}(B\mid y)=\frac{\Pi^{\otimes}(B\cap\mathcal C\mid y)}{\Pi^{\otimes}(\mathcal C\mid y)} ,
\]
so the joint posterior is the product of the individual posteriors conditioned on $\mathcal C$. The proof has three parts: the conditioning costs at most the factor $1/\Pi^{\otimes}(\mathcal C\mid y)$ (Lemma~\ref{pf:lem:cond}); the true curves are separated (Lemma~\ref{pf:lem:gap}); and each individual posterior concentrates in the supremum norm (Lemma~\ref{pf:lem:sup}), so that $\Pi^{\otimes}(\mathcal C\mid y)\ge1/2$ with high probability. Theorem~\ref{thm:main} is then applied at each level. We take $\sigma_j=1$ for all $j$ until the last paragraph of the proof. In this subsection $J$ is the number of levels, not the number of knots of Proposition~\ref{pf:prop:shell}.

\begin{lemma}[conditioning]\label{pf:lem:cond}
Let $n>r$. For every $y$, $\Pi^{\otimes}(\mathcal C\mid y)>0$, so $\Pi_{\mathcal C}(\cdot\mid y)$ is a probability measure. For every $j$ and every Borel set $A\subset\R^n$,
\[
\Pi_{\mathcal C}(\Theta:\ \theta_j\in A\mid y)\le\frac{\Pi_j(A\mid y)}{\Pi^{\otimes}(\mathcal C\mid y)} .
\]
\end{lemma}

\begin{proof}
By Lemma~\ref{pf:lem:coord}(b) the normalizing integral in \eqref{pf:eq:post} is finite, so each $\Pi_j(\cdot\mid y)$ has a Lebesgue density that is positive at every point of $\R^n$, and $\Pi^{\otimes}(\cdot\mid y)$ has a density that is positive at every point of $(\R^n)^J$. The set $\mathcal C$ contains the nonempty open set $\{\Theta:\ \theta_{i1}<\cdots<\theta_{iJ},\ i=1,\dots,n\}$, which has positive Lebesgue measure. Hence $\Pi^{\otimes}(\mathcal C\mid y)>0$. For the inequality, $\Pi^{\otimes}(\{\theta_j\in A\}\cap\mathcal C\mid y)\le\Pi^{\otimes}(\theta_j\in A\mid y)=\Pi_j(A\mid y)$.
\end{proof}

\begin{lemma}[separation]\label{pf:lem:gap}
Let Assumption~\ref{ass:B} hold at the levels $\tau_1<\cdots<\tau_J$ with the constants $L>0$ and $0<\fu<\infty$, and let
\[
g_0=\min\Bigl\{L,\ \min_{1\le j<J}\frac{\tau_{j+1}-\tau_j}{\fu}\Bigr\}.
\]
Then $\thetas_{i,j+1}-\thetas_{ij}\ge g_0$ for all $n$, all $i$ and all $1\le j<J$.
\end{lemma}

\begin{proof}
Fix $i$ and $j<J$, and write $F=F_{y_i}$. Assumption~\ref{ass:B} at the level $\tau_j$ makes $F$ continuous at the $\tau_j$-quantile $\thetas_{ij}$; since $F(\thetas_{ij}-)\le\tau_j\le F(\thetas_{ij})$, this gives $F(\thetas_{ij})=\tau_j$. Also $F(\thetas_{i,j+1})\ge\tau_{j+1}>\tau_j$ because $\thetas_{i,j+1}$ is a $\tau_{j+1}$-quantile. Since $F$ is nondecreasing, $d=\thetas_{i,j+1}-\thetas_{ij}>0$. If $d\le L$, Assumption~\ref{ass:B} at the level $\tau_j$ gives
\[
\tau_{j+1}-\tau_j\le F(\thetas_{ij}+d)-F(\thetas_{ij})\le\fu\,d .
\]
So either $d>L$ or $d\ge(\tau_{j+1}-\tau_j)/\fu$, and in both cases $d\ge g_0$.
\end{proof}

The next lemma concerns one level. It is stated in the setting and with the notation of Theorem~\ref{thm:main} and of its proof: $\rho_n$, $T_n$ and $K_n$ are defined before Proposition~\ref{pf:prop:local}, $C_I$ is the constant of Lemma~\ref{pf:lem:interp}, and $\Omega_E$, $\Omega_G$ and $n_2$ are those of Lemma~\ref{pf:lem:evid} and Proposition~\ref{pf:prop:local}.

\begin{lemma}[supremum-norm concentration]\label{pf:lem:sup}
In the setting of Theorem~\ref{thm:main} with $\sigma=1$, let
\[
s_n=C_I\bigl\{\rho_n+(\rho_n^{2r-2}T_n)^{1/(2r-1)}\bigr\},\qquad \eta_n=r!\,\bigl\{e^{-\fl n\rho_n^2/16}+e^{-n/8}\bigr\}.
\]
Then $s_n\to0$ and $\eta_n\to0$; the sequence $s_n$ depends only on $(c,r)$ and $\eta_n$ only on $(\fl,r)$. For $n\ge n_2$,
\[
\Pi_\lambda\bigl(\theta:\ \|\theta-\thetas\|_\infty>s_n\bigm|y\bigr)\le\eta_n\quad\text{on }\Omega_E\cap\Omega_G,
\qquad
\pr\bigl\{(\Omega_E\cap\Omega_G)^c\bigr\}\le\Bigl(\frac54+\frac{64}{\fl^2}\Bigr)n^{-1/(2r)} .
\]
\end{lemma}

\begin{proof}
For $\delta\in K_n$, Lemma~\ref{pf:lem:interp} with $\nn{\delta}\le\rho_n$ and $\TV^{(r)}(\delta)\le T_n$ gives $\|\delta\|_\infty\le s_n$. Hence $\{\theta:\ \|\theta-\thetas\|_\infty>s_n\}\subset\{\theta:\ \theta-\thetas\notin K_n\}$, and Proposition~\ref{pf:prop:local}(b) bounds the posterior mass of the larger set by $\eta_n$ on $\Omega_E\cap\Omega_G$. Since $\rho_n^{2r-2}T_n=(4/c)\,n^{-(4r^2-8r+1)/(2(2r+1))}(\log n)^{r-1}\to0$ for $r\ge2$, as in the proof of Proposition~\ref{pf:prop:local}(a), and $\rho_n\to0$, we get $s_n\to0$; since $n\rho_n^2\to\infty$, $\eta_n\to0$. For the probability, Lemma~\ref{pf:lem:evid} and Proposition~\ref{pf:prop:local} give $\pr(\Omega_E^c)+\pr(\Omega_G^c)\le\tfrac54(n\en^2)^{-1}+64/(\fl^2n\rho_n^2)$, and $n\rho_n^2=n^{2/(2r+1)}\log n\ge n^{1/(2r)}=n\en^2$, because $2/(2r+1)\ge1/(2r)$ and $\log n\ge1$ for $n>r\ge2$.
\end{proof}

\begin{proof}[Proof of Theorem~\ref{thm:joint}]
For each level $l$, the hypotheses of Theorem~\ref{thm:main} hold at $\tau=\tau_l$ with $c=c_l$. Let $M^{(l)}$, $C^{(l)}$ and $n_0^{(l)}$ be the constants of that theorem, and let $s_{n,l}$, $\Omega_l=\Omega_{E,l}\cap\Omega_{G,l}$ and $n_2^{(l)}$ be the sequence, the event and the threshold of Lemma~\ref{pf:lem:sup} at that level; $\eta_n$ is the same at all levels. Put
\[
M=\max_{l}M^{(l)},\qquad C=2\max_{l}C^{(l)}+J\Bigl(\frac54+\frac{64}{\fl^2}\Bigr),\qquad \Omega=\bigcap_{l=1}^J\Omega_l ,
\]
and let $n_0\ge\max_l\max\{n_0^{(l)},n_2^{(l)}\}$ be such that $\max_ls_{n,l}\le g_0/2$ and $J\eta_n\le1/2$ for all $n\ge n_0$, with $g_0$ as in Lemma~\ref{pf:lem:gap}. Such an $n_0$ exists by Lemma~\ref{pf:lem:sup}, and it depends only on the quantities listed in the theorem. Let $n\ge n_0$.

\emph{The constraint.} If $\|\theta_l-\thetas_l\|_\infty\le g_0/2$ for every $l$, then by Lemma~\ref{pf:lem:gap}, for all $i$ and $j<J$,
\[
\theta_{i,j+1}-\theta_{ij}\ge(\thetas_{i,j+1}-\thetas_{ij})-g_0\ge0 ,
\]
that is, $\Theta\in\mathcal C$. Hence, since $s_{n,l}\le g_0/2$, a union bound and Lemma~\ref{pf:lem:sup} give, on $\Omega$,
\[
\Pi^{\otimes}(\mathcal C^c\mid y)\le\sum_{l=1}^J\Pi_l\bigl(\theta_l:\ \|\theta_l-\thetas_l\|_\infty>s_{n,l}\bigm|y\bigr)\le J\eta_n\le\frac12 .
\]

\emph{Transfer.} Fix $j$ and let $A=\{\theta\in\R^n:\ \Delta_n^2(\theta-\thetas_j)>M\en^2\}$. By Lemma~\ref{pf:lem:cond} and the last display, $\Pi_{\mathcal C}(\theta_j\in A\mid y)\le2\,\Pi_j(A\mid y)$ on $\Omega$, and $\Pi_{\mathcal C}\le1$ on $\Omega^c$. Since $M\ge M^{(j)}$, Theorem~\ref{thm:main} at the level $\tau_j$ bounds $\E_{\thetas}\Pi_j(A\mid y)$ by $C^{(j)}n^{-1/(2r)}$. With the probability bound of Lemma~\ref{pf:lem:sup} at each level,
\[
\begin{aligned}
\E_{\thetas}\,\Pi_{\mathcal C}(\theta_j\in A\mid y)&\le\pr(\Omega^c)+2\,\E_{\thetas}\Pi_j(A\mid y)\\
&\le J\Bigl(\frac54+\frac{64}{\fl^2}\Bigr)n^{-1/(2r)}+2C^{(j)}\,n^{-1/(2r)}\le C\,n^{-1/(2r)} .
\end{aligned}
\]
All constants are those of Theorem~\ref{thm:main} and of Lemma~\ref{pf:lem:sup}, so the bound is uniform over the class. The constant $M$ is the largest of the $J$ single-level constants, and neither $M$ nor $C$ depends on the spacing of the levels; the spacing enters $n_0$ only, through $g_0$.

\emph{General scales.} Let $\sigma_j=\sigma_{0j}>0$. Lemmas~\ref{pf:lem:cond} and~\ref{pf:lem:gap} do not involve the scales. By the reduction to $\sigma=1$ at the beginning of this appendix, the law of $\theta_l/\sigma_{0l}$ under $\Pi_l(\cdot\mid y)$ is the $\sigma=1$ posterior based on $y/\sigma_{0l}$ with $c_l$ replaced by $\sigma_{0l}c_l$, and the rescaled data satisfy the assumptions with $(L/\sigma_{0l},\sigma_{0l}\fl,\sigma_{0l}\fu)$ and $\Vb/\sigma_{0l}$. Lemma~\ref{pf:lem:sup} for the rescaled problem gives a sequence $s_{n,l}'\to0$, a sequence $\eta_{n,l}'\to0$ and an event $\Omega_l$ with $\pr(\Omega_l^c)\le\{5/4+64/(\sigma_{0l}\fl)^2\}n^{-1/(2r)}$, on which $\Pi_l(\|\theta_l-\thetas_l\|_\infty>\sigma_{0l}s_{n,l}'\mid y)\le\eta_{n,l}'$. The two steps above apply with $\sigma_{0l}s_{n,l}'$ in place of $s_{n,l}$, $\sum_l\eta_{n,l}'$ in place of $J\eta_n$, and Theorem~\ref{thm:main} for the scale $\sigma_{0j}$. This gives the bound with $M=\max_lM^{(l)}$ and $C=2\max_lC^{(l)}+\sum_l\{5/4+64/(\sigma_{0l}\fl)^2\}$, where $M^{(l)}$, $C^{(l)}$ and $n_0^{(l)}$ are now the constants of Theorem~\ref{thm:main} at the scale $\sigma_{0l}$ (in the original units, so $M^{(l)}$ includes the factor $\max\{\sigma_{0l},\sigma_{0l}^2\}$ of the reduction), and with $n_0$ at least every $n_0^{(l)}$ and every threshold $n_2$ of the rescaled problems, and such that $\max_l\sigma_{0l}s_{n,l}'\le g_0/2$ and $\sum_l\eta_{n,l}'\le1/2$ for $n\ge n_0$.
\end{proof}

The proof gives slightly more than is stated. First, for $n\ge n_0$, on $\Omega$ and for any Borel sets $A_1,\dots,A_J\subset\R^n$, $\Pi_{\mathcal C}(\cup_j\{\theta_j\in A_j\}\mid y)\le2\sum_j\Pi_j(A_j\mid y)$. With $A_j=\{\theta:\ \Delta_n^2(\theta-\thetas_j)>M\en^2\}$, the bound of Theorem~\ref{thm:joint} therefore holds simultaneously for all levels, that is, for $\max_j\Delta_n^2(\theta_j-\thetas_j)$, with $C$ replaced by $2\sum_lC^{(l)}+J(5/4+64/\fl^2)$ at unit scales. Second, $\sup_B|\Pi_{\mathcal C}(B\mid y)-\Pi^{\otimes}(B\mid y)|=\Pi^{\otimes}(\mathcal C^c\mid y)$. For $n\ge n_0$, outside an event of probability at most $n^{-1/(2r)}\sum_l\{5/4+64/(\sigma_{0l}\fl)^2\}$, the ordering constraint therefore changes the posterior by at most $\sum_l\eta_{n,l}'$ in this distance; at unit scales these are $J(5/4+64/\fl^2)n^{-1/(2r)}$ and $J\eta_n$. The distance bound $\sum_l\eta_{n,l}'$ tends to zero faster than any negative power of $n$; the probability of the exceptional event is at most of order $n^{-1/(2r)}$.


\end{document}
